\documentclass[11pt,a4paper]{article}
\usepackage[margin=2.35cm]{geometry}
\usepackage[T1]{fontenc}
\usepackage[utf8]{inputenc}
\usepackage{lmodern}
\usepackage{microtype}
\usepackage{amsmath,amssymb,mathtools}
\usepackage{booktabs}
\usepackage{graphicx}
\usepackage{hyperref}
\usepackage{xcolor}
\usepackage{enumitem}
\usepackage{siunitx}
\usepackage{array}
\usepackage{fancyhdr}
\usepackage{titlesec}
\usepackage{abstract}
\usepackage{url}
\usepackage{physics}
\hypersetup{colorlinks=true,linkcolor=black,citecolor=blue!55!black,urlcolor=blue!55!black,pdftitle={From Exact Causal Binding to a Falsifiable Planck-Tick Gap Law},pdfauthor={Ingolf Lohmann}}
\setlength{\parindent}{0pt}
\setlength{\parskip}{0.55em}
\pagestyle{fancy}\fancyhf{}\lhead{QIK-VRT Planck-Tick Gap Law v1}\rhead{3 September 2026}\cfoot{\thepage}
\newcommand{\tP}{t_{\mathrm P}}\newcommand{\EP}{E_{\mathrm P}}\newcommand{\dE}{\Delta E}\newcommand{\dw}{\Delta\omega}\newcommand{\repo}{\texttt{Goldkelch/qik-vrt}}
\title{\textbf{From Exact Causal Binding to a Falsifiable Planck-Tick Gap Law}\\[0.35em]\large A QIK-VRT physical hypothesis with a parameter-free phase-frequency residual}
\author{\textbf{Ingolf Lohmann}\\\small Strategic conception and authorship\\[0.2em]\small Mathematical derivation and technical preparation with OpenAI}
\date{3 September 2026}
\begin{document}\maketitle
\begin{center}\fbox{\parbox{0.93\textwidth}{\small\textbf{Scientific status.} This manuscript freezes a falsifiable physical hypothesis and a differentiating numerical prediction. It does \emph{not} claim empirical confirmation, independent reproduction, a detected Planck-time clock, or final physical closure. Repository lineage: \repo, base \texttt{main@03759aebbd983602b3ddfe7bf751a79535559a6f}, physical-hypothesis carrier PR \#962.}}\end{center}
\begin{abstract}QIK-VRT was developed as an exact-binding and causal-evidence framework in which claims are attached to exact system states, transport acknowledgement is separated from effect acknowledgement, and a changed state requires reobservation rather than transfer of predecessor evidence. Earlier repository artifacts formalized symbolic Planck identities, identity-preserving wave/record projections, measurement-derived dimensions, and fail-closed boundaries between formal proof and physical discovery. The remaining physics threshold was a differentiating observable. This paper crosses that threshold by freezing a concrete physical postulate: for a bound quantum system, the ground-referenced Hamiltonian $H_c=H-E_0I$ obeys a symmetric finite-difference evolution law with universal time step equal to the Planck time $\tP=\sqrt{\hbar G/c^5}$. For an excitation gap $\dE=E-E_0$, the low-energy branch predicts $\omega_Q=\arcsin(\dE/\EP)/\tP$, with $\EP=\hbar/\tP$, instead of $\omega_0=\dE/\hbar$. The exact residual $\dw=[\arcsin(\dE/\EP)-\dE/\EP]/\tP>0$ is parameter-free once $\hbar$, $G$, and $c$ are fixed. Its low-energy leading term is cubic in the excitation gap. This makes the proposal experimentally falsifiable and distinguishes it from a mere reparameterization.\end{abstract}
\section{Repository and scientific lineage}The present hypothesis is deliberately continuous with earlier QIK-VRT artifacts rather than a disconnected conjecture. The repository already contains exact symbolic Planck normal forms, identity-preserving observation, measurement-derived dimensions, and a fail-closed epistemic boundary. The present work supplies the previously missing ingredient: a specific dynamics with a quantitative observable that differs from ordinary continuous-time quantum mechanics.
\section{Why a physical threshold requires a non-reparameterization theorem}Let $R:X\to X'$ be an invertible representation map, $T:X\to X$ the original dynamics, and $O:X\to Y$ an observable. With $T'=RTR^{-1}$ and $O'=OR^{-1}$, $O'(T'(Rx))=O(Tx)$. Thus an invertible reparameterization cannot by itself generate a differentiating observable. A physical hypothesis requires a dynamics $T_Q$ not reducible to $RTR^{-1}$ and at least one observable with $O(T_Qx)\neq O(Tx)$.
\section{Physical postulate PTG-P01}For a bound quantum system with self-adjoint Hamiltonian $H$ and ground energy $E_0$, define $H_c=H-E_0I$. Define $\tP=\sqrt{\hbar G/c^5}$ and $\EP=\hbar/\tP$. \textbf{PTG-P01:} $i\hbar[\psi(t+\tP)-\psi(t-\tP)]/(2\tP)=H_c\psi(t)$. This is a physical postulate, not a theorem of standard quantum mechanics.
\section{Exact gap law}For an energy eigenstate with $\dE=E-E_0$ and $\psi(t)=e^{-i\omega t}\psi(0)$, substitution gives $\hbar\sin(\omega\tP)/\tP=\dE$. With $x=\dE/\EP$ and the low-energy branch $0\le x<1$, \[\boxed{\omega_Q=\frac{\arcsin(\dE/\EP)}{\tP}}.\] Standard quantum mechanics gives $\omega_0=\dE/\hbar=x/\tP$, hence \[\boxed{\dw=\frac{\arcsin(x)-x}{\tP}>0\quad(0<x<1).}\]
\section{Low-energy expansion and scaling law}Using $\arcsin x=x+x^3/6+3x^5/40+\cdots$, \[\boxed{\dw=\frac{\dE^3}{6\hbar\EP^2}+O\!\left(\frac{\dE^5}{\hbar\EP^4}\right)}\] and \[\boxed{\frac{\omega_Q}{\omega_0}-1=\frac16\left(\frac{\dE}{\EP}\right)^2+O\!\left((\dE/\EP)^4\right).}\]
\section{Parameter count and falsifiability}In v1 there is no deformation coefficient. For a coherently phase-tracked bound two-level system with independently determined $\dE$, define $R=\omega_{\mathrm{measured}}-\dE/\hbar$. PTG-P01 predicts $R_Q=[\arcsin(\dE/\EP)-\dE/\EP]/\tP>0$. A validated bound $|R|<R_Q$, with total uncertainty and systematic error below $R_Q$, falsifies PTG-P01 in that domain. Agreement at one gap is insufficient; distinct-gap measurements must reproduce the predicted cubic low-energy scaling.
\section{Numerical predictions}Using CODATA-2022 values gives $\tP\approx5.3912464467\times10^{-44}$ s and $\EP\approx1.2208901282\times10^{28}$ eV. Frozen examples: at 1 eV, relative shift $1.1181\times10^{-57}$ and $\dw=1.6988\times10^{-42}$ rad/s; at 1 MeV, $1.1181\times10^{-45}$ and $1.6988\times10^{-24}$ rad/s; at 1 GeV, $1.1181\times10^{-39}$ and $1.6988\times10^{-15}$ rad/s; at 1 TeV, $1.1181\times10^{-33}$ and $1.6988\times10^{-6}$ rad/s. These are prediction targets, not claims of present experimental feasibility.
\section{Connection to QIK-VRT exact causal binding}The prediction is bound to a declared Hamiltonian, ground reference, excitation gap, and branch. Measurements of other transitions do not inherit predecessor disposition. Preparing or transmitting a control sequence is not an effect acknowledgement; phase must be read back. Changes in Hamiltonian, calibration, environment, or branch require reobservation. Mathematical derivation does not imply that nature realizes the postulate.
\section{Relation to prior finite-difference and chronon work}Finite-difference time in quantum mechanics is historically older. Caldirola and later Farias/Recami considered chronon formulations, including symmetric finite differences. No priority claim is made for discretizing the Schroedinger equation. The narrower v1 commitment is the universal Planck-time tick, ground-referenced binding, selected symmetric low-energy branch, exact arcsin gap law, and absence of a free deformation coefficient.
\section{Domain and cautions}The equation has branch structure because $\sin(\omega\tP)=x$; v1 explicitly selects the low-energy branch continuously connected to standard quantum mechanics. The domain $0\le\dE<\EP$ is declared. Relativistic covariance, quantum-field-theoretic completion, arbitrary-system unitarity, Standard Model limit, and Einstein limit remain separate bridge obligations.
\section{Experimental program}A valid test requires an exactly specified bound two-level or effective two-level system, independent determination of $\dE$, phase-frequency readout with a complete uncertainty budget, preregistered raw-data mapping, distinct-gap testing of cubic scaling, and independent reproduction. Ordinary atomic and optical gaps imply extremely small residuals; metrology, long coherent accumulation, interferometric amplification, composite effective gaps, and high-energy coherent systems are natural feasibility targets, but no amplification scheme may introduce an unfrozen fitting parameter.
\section{What has and has not been achieved}The formalism-to-physics threshold is crossed in the limited sense that QIK-VRT now makes a concrete physical assertion that can be false. Physical hypothesis bound: true. Differentiating prediction frozen: true. Experimental platform bound: false. Measurement performed: false. Empirical correspondence: false. Independent reproduction: false. PASS, FINAL\_PASS and EFFECT\_ACK\_DONE: false.
\section{Repository reproducibility map}The canonical machine-readable v1 contract is \texttt{physics/planck\_tick\_gap\_law\_v1.json}; the human derivation is \texttt{docs/research/2026-09-03-planck-tick-gap-law/README.md}; the carrier is PR \#962. Antecedent formal artifacts include \texttt{VRTCore\_SMG\_PlanckBridge.lean}, \texttt{VRTCore\_SMG\_AxiomAudit.lean}, and \texttt{MeasurementDerivedDimensions.lean}. A broader prior QIK-VRT formalization is indexed on Zenodo as DOI \href{https://doi.org/10.5281/zenodo.21488116}{10.5281/zenodo.21488116}.
\section{Conclusion}Under PTG-P01, a bound excitation gap obeys the parameter-free low-energy phase law $\omega_Q=\arcsin(\dE/\EP)/\tP$. The difference from standard quantum mechanics is positive, quantitative, and asymptotically cubic in the gap. The repository can prove the derivation, bind constants, freeze the prediction, and enforce the evidence boundary. Only observation can determine whether nature realizes the postulate.
\section*{References}\begin{enumerate}\item P. J. Mohr, D. B. Newell, B. N. Taylor, and E. Tiesinga, ``CODATA recommended values of the fundamental physical constants: 2022,'' \emph{Rev. Mod. Phys.} 97, 025002 (2025), DOI 10.1103/RevModPhys.97.025002.\item R. A. H. Farias and E. Recami, ``Introduction of a Quantum of Time (chronon), and its Consequences for Quantum Mechanics,'' arXiv:quant-ph/9706059 (1997).\item P. Caldirola, chronon / finite-difference dynamics program; see Ref. 2 for review and bibliography.\item I. Lohmann, QIK-VRT repository, \url{https://github.com/Goldkelch/qik-vrt}.\item I. Lohmann, QIK-VRT formalization, Zenodo DOI 10.5281/zenodo.21488116.\end{enumerate}
\end{document}
